\section{\texorpdfstring{Dihedral Quotients of the Tower and the Field $E_{W''}$}{Dihedral Quotients of the Tower and the Field E\_W''}}
\label{dh:section}

Every field $K\in\mathcal K$ contains the fixed field of $D_3G_B$
(Theorem~\ref{tw:field-family}), whose Galois group over $B$ is the extension
$G_B/D_3G_B$ of $G_B/D_2G_B=\Gal(E/B)=\Ftwo^8$ by
$\operatorname{gr}_2G_B=D_2G_B/D_3G_B$, of dimension $15$
(Lemma~\ref{tw:quadratic-layer}). The Kummer field $E$ sees only the quotient
$\Ftwo^8$. In particular, a prime of $B$ whose Frobenius vector is $0$ splits
completely in $E$, and the comparison with $E$ treats it as if it split
completely in $K$. This section uses the layer $\operatorname{gr}_2G_B$:
through dihedral quotients of $G_B$ it gives lower bounds for residue
degrees, by a census for primes of norm at most $4\cdot10^{13}$ and through
the field $\EW$ for all norms, and it supplies the system of floors of
Definition~\ref{dh:floors}.

Subsection~\ref{dh:forms-section} identifies the linear forms on
$\operatorname{gr}_2G_B$ with quadratic forms and finds the $27$ forms of
dihedral type; Subsection~\ref{dh:fields-section} realizes each of them by a
dihedral field contained in every $K\in\mathcal K$.
Subsection~\ref{dh:W-section} defines the space $W''$ spanned by four of
these forms. Subsection~\ref{dh:census-section} reports the census of the
primes of degree one with Frobenius vector $0$.
Subsection~\ref{dh:EW-section} studies the field $\EW$ and factors its zeta
function into the $L$-functions evaluated in Section~\ref{lv:section}, and
Subsection~\ref{dh:floors-section} defines the floors and expresses their
floor sum through $\zeta_{\EW}$, for use in Section~\ref{ce:section}.

In Sections~\ref{dh:section}--\ref{ce:section}, supplementary programs and
data files are named by their paths relative to the directory
\texttt{certificates/dihedral} if they lie in it, and relative to
\texttt{certificates} otherwise; the names are unique.

For $r,v\in\Ftwo^8$ we write $r\cdot v=\sum_ir_iv_i\in\Ftwo$, and, as in
Definition~\ref{gf:characters}, $\alpha_r=\prod_i\alpha_i^{r_i}$ and
$\chi_r(g)=(-1)^{r\cdot\bar g}$, where $\bar g\in\Ftwo^8$ is the elementary
image of $g$ (Definition~\ref{tw:elementary-image}). We use the same symbol
$\chi_r$ for the inflation of this character to any Galois group over $B$
that maps onto $\Gal(E/B)$. We write $\mathrm{Dih}_4$ for the dihedral group
of order $8$.

\subsection{\texorpdfstring{Linear forms on $\operatorname{gr}_2G_B$}{Linear forms on gr\_2 G\_B}}
\label{dh:forms-section}
This subsection identifies the linear forms on $\operatorname{gr}_2G_B$
with quadratic forms on $\Ftwo^8$ (Lemma~\ref{dh:quadratic}) and isolates the
$27$ forms whose polar form has rank $2$ (Lemma~\ref{dh:27}). All of them
are of dihedral type (Definition~\ref{dh:dihedral-type}), and the next
subsection realizes each of them by a dihedral field.

Recall (Lemma~\ref{tw:graded}) that $\operatorname{gr}_2G_B=\mathfrak
L_2/R_2$, where $\mathfrak L_2$ is the degree-two part of the free
restricted Lie algebra on $\mathfrak L_1=\Ftwo^8$ and $R_2$ is the relation
space (Definition~\ref{tw:relation-space}). A \emph{quadratic form} on
$\Ftwo^8$ is a function $Q\colon\Ftwo^8\to\Ftwo$ such that
$Q(v+w)-Q(v)-Q(w)$ is bilinear in $v,w$; this bilinear form is its
\emph{polar form}. It is alternating, and its rank is even.

\begin{lemma}\label{dh:quadratic}
For a linear form $\psi$ on $\mathfrak L_2$ put $q_\psi(v)=\psi(v^{[2]})$.
Then $\psi\mapsto q_\psi$ is a linear bijection from the dual of
$\mathfrak L_2$ onto the space of quadratic forms on $\Ftwo^8$, and the polar
form of $q_\psi$ is $\operatorname{pol}_\psi(v,w)=\psi([v,w])$. The linear
forms on $\operatorname{gr}_2G_B$ are the $\psi$ that vanish on $R_2$.
\end{lemma}

\begin{proof}
The identity $(v+w)^{[2]}=v^{[2]}+w^{[2]}+[v,w]$ shows that $q_\psi$ is a
quadratic form with polar form $\psi([v,w])$. The map is injective, because
$\mathfrak L_2$ is spanned by the restricted squares $v^{[2]}$, the brackets
being differences of them. Both spaces have dimension $36$.
\end{proof}

\begin{definition}\label{dh:dihedral-type}
A linear form $\psi$ on $\operatorname{gr}_2G_B$ is \emph{of dihedral type} if
there are linearly independent $r_1,r_2\in\Ftwo^8$ and
$\mathbf r\in\Ftwo^2\setminus\{0\}$ such that, with
$\phi(v)=(r_1\cdot v,r_2\cdot v)$,
\[
 q_\psi(v)=\begin{cases}1&\text{if }\phi(v)=\mathbf r,\\0&\text{otherwise.}
 \end{cases}
\]
We call $(r_1,r_2,\mathbf r)$ \emph{data} for $\psi$, and $\phi$ its
\emph{associated map} $\Ftwo^8\to\Ftwo^2$.
\end{definition}

\begin{lemma}\label{dh:dihedral-polar}
Let $\psi$ be of dihedral type with data $(r_1,r_2,\mathbf r)$ and associated
map $\phi$. Then $\operatorname{pol}_\psi(v,w)=\det(\phi(v),\phi(w))$; it has
rank $2$, and its radical is $\ker\phi$.
\end{lemma}

\begin{proof}
The function $x\mapsto[x=\mathbf r]$ on $\Ftwo^2$ is $x_1x_2$, $x_1+x_1x_2$
or $x_2+x_1x_2$ for $\mathbf r=(1,1),(1,0),(0,1)$. The linear terms have
polar form $0$, and $x_1x_2$ has polar form $x_1y_2+x_2y_1=\det(x,y)$. Since
$q_\psi(v)=[\phi(v)=\mathbf r]$ and $\phi$ is linear,
$\operatorname{pol}_\psi(v,w)=\det(\phi(v),\phi(w))$. As $r_1,r_2$ are
linearly independent, $\phi$ is surjective, and $\det$ is nondegenerate on
$\Ftwo^2$; so the radical is $\ker\phi$, of codimension $2$.
\end{proof}

\begin{lemma}\label{dh:27}
Exactly $27$ of the $2^{15}-1$ nonzero linear forms on $\operatorname{gr}_2G_B$
have a polar form (Lemma~\ref{dh:quadratic}) of rank $2$. All of them are of
dihedral type (Definition~\ref{dh:dihedral-type}), and they span the dual of
$\operatorname{gr}_2G_B$. The polar forms of the other nonzero forms have rank
$4$, $6$ or $8$. We number the $27$ forms $0$ to $26$ as in the file
\texttt{d4all27.json}, and write $\psi_o$ for the form number $o$.
\end{lemma}

\begin{proof}[Proof (computer-assisted)]
The supplementary programs \texttt{d4search.py} and \texttt{plane.py}
enumerate the linear forms on $\mathfrak L_2$ that vanish on the space $R_2$
of \texttt{lie241.py} (Lemma~\ref{tw:quadratic-layer}), compute the ranks of
their polar forms, and, for those of rank $2$, find data $(r_1,r_2,\mathbf r)$;
the $27$ forms and their data are listed in \texttt{d4all27.json}.
\end{proof}

\subsection{Dihedral fields}\label{dh:fields-section}
This subsection realizes each form of dihedral type by a Galois extension
$M/B$ with group $\mathrm{Dih}_4$, contained in every $K\in\mathcal K$, and
shows how $M$ forces even residue degrees (Lemma~\ref{dh:detection}); the
elements that define these fields exist for all $27$ forms
(Lemma~\ref{dh:beta-exists}). The census and the field $\EW$ are built from
these fields.

\begin{lemma}\label{dh:detection}
Let $\psi$ be a form of dihedral type with data $(r_1,r_2,\mathbf r)$
(Definition~\ref{dh:dihedral-type}). Put
$F_0=B(\sqrt{\alpha_{r_1}},\sqrt{\alpha_{r_2}})$, a biquadratic extension of
$B$ contained in $E$, and identify $\Gal(F_0/B)$ with $\Ftwo^2$ through the
signs by which an element multiplies $\sqrt{\alpha_{r_1}}$ and
$\sqrt{\alpha_{r_2}}$. Let $\beta\in F_0^\times$ be a unit outside the primes
above $S$ such that, for each $\sigma\ne1$ in $\Gal(F_0/B)$ (in this lemma
$\sigma$ denotes such an element):
\begin{enumerate}
\item[(D1)] $\sigma(\beta)/\beta=u_\sigma^2$ for some $u_\sigma\in F_0^\times$;
\item[(D2)] the number $t_\sigma=u_\sigma\sigma(u_\sigma)$, which is $\pm1$, is
 $-1$ exactly when $\sigma=\mathbf r$.
\end{enumerate}
Let $M=F_0(\sqrt\beta)$. Then:
\begin{enumerate}
\item[(a)] $M$ is Galois over $B$ with group isomorphic to $\mathrm{Dih}_4$.
 The elements of $\Gal(M/B)$ with image $\mathbf r$ in $\Gal(F_0/B)$ have
 order $4$.
\item[(b)] $M$ lies in the fixed field of $D_3G_B$, hence in every
 $K\in\mathcal K$. The quotient map $\pi_M\colon G_B\to\Gal(M/B)$ maps
 $D_2G_B$ to the center $\langle z\rangle$ of $\Gal(M/B)$, by
 $g\mapsto z^{\psi([g]_2)}$, where $[g]_2$ is the class of $g$ in
 $\operatorname{gr}_2G_B$.
\item[(c)] Let $\mathfrak p\notin S$ be a prime of $B$ that splits completely
 in $F_0$, and suppose that $\beta$ is not a square in the residue field of a
 prime of $F_0$ above $\mathfrak p$. Then $f_K(\mathfrak p)$ is even for every
 $K\in\mathcal K$. The hypothesis does not depend on the prime of $F_0$ above
 $\mathfrak p$.
\item[(d)] Let $\mathfrak p\notin S$ have Frobenius vector $0$. Then the
 Frobenius elements in $G_B$ of the primes of $B_G$ above $\mathfrak p$ lie in
 $D_2G_B$, and they have the same class
 $\overline{\Frob}^{(2)}_{\mathfrak p}\in\operatorname{gr}_2G_B$. The
 hypothesis of (c) holds exactly when
 $\psi(\overline{\Frob}^{(2)}_{\mathfrak p})=1$.
\end{enumerate}
\end{lemma}

\begin{proof}
Since $\sigma^2=1$, $t_\sigma^2=(\sigma(\beta)/\beta)\cdot\sigma(\sigma(\beta)/\beta)
=\sigma^2(\beta)/\beta=1$, so $t_\sigma=\pm1$.

(a) Every conjugate of $\beta$ over $B$ is $\beta$ times a square in $F_0$, so
$M$ is Galois over $B$. If $\beta$ were a square in $F_0$, we could take
$u_\sigma=\sigma(\sqrt\beta)/\sqrt\beta$, and then $t_\sigma=1$ for every
$\sigma$, contrary to (D2); so $[M:B]=8$. Let $\tilde\sigma\in\Gal(M/B)$ lift
$\sigma\ne1$. Then $\tilde\sigma(\sqrt\beta)=\eta u_\sigma\sqrt\beta$ with
$\eta=\pm1$, and
$\tilde\sigma^2(\sqrt\beta)=\eta\sigma(u_\sigma)\,\eta u_\sigma\sqrt\beta
=t_\sigma\sqrt\beta$. So $\tilde\sigma^2$ is the generator $z$ of
$\Gal(M/F_0)$ exactly when $\sigma=\mathbf r$. The group $\Gal(M/B)$ has order
$8$, a central subgroup $\langle z\rangle$ (a normal subgroup of order two),
and quotient $\Ftwo^2$; among such groups, $C_2^3$ has no element of order
four, $C_4\times C_2$ has elements of order four above two of the three
nonzero classes of the quotient, the quaternion group above all three, and
$\mathrm{Dih}_4$ above exactly one. Hence $\Gal(M/B)\cong\mathrm{Dih}_4$. In
$\mathrm{Dih}_4$, an element squares to $z$ exactly when its image in
$\Ftwo^2$ is $\mathbf r$, and two elements have commutator $z$ exactly when
their images are linearly independent.

(b) The field $F_0$ lies in $E$, so it is unramified over $B$ outside $S$,
and $\beta$ is a unit outside the primes above $S$, which include all primes
above $2$; so $M/F_0$ is also unramified outside $S$. Since $[M:B]=8$, $M$
lies in $B_S$, and restriction gives a surjection $\pi_M\colon
G_S\to\Gal(M/B)$. If $g\in G_S$ has elementary image $v$, then $\pi_M(g)$ has
image $\phi(v)$ in $\Gal(F_0/B)$, because $g(\sqrt{\alpha_r})=(-1)^{r\cdot v}
\sqrt{\alpha_r}$. By (a) and Lemma~\ref{dh:dihedral-polar}, for $g,h\in G_S$
with elementary images $v,w$,
\[
 \pi_M(g^2)=z^{q_\psi(v)},\qquad \pi_M([g,h])=z^{\operatorname{pol}_\psi(v,w)}.
\]
We check that $\pi_M$ kills the generators (ii)--(iv) of $N_B$ in
Lemma~\ref{tw:GB-presentation}, more precisely their images in $G_S$, which
generate $N$ as a normal subgroup. The elements $x_j^2$, $[x_j,y_j]$,
$[x_j,z_j]$, $\tau_\nu^2$ and $\varphi_\nu^2$ map to $z^{\psi(\xi)}$, where
$\xi$ is respectively $\bar x_j^{\,[2]}$, $[\bar x_j,\bar y_j]$,
$[\bar x_j,\bar z_j]$, $\bar\tau_\nu^{\,[2]}$ and $\bar\varphi_\nu^{\,[2]}$;
these lie in $R_2$ by \eqref{tw:initials-list}, and $\psi$ vanishes on $R_2$.
The elements $[[y_j,z_j],z_j]$, $z_j^4$, $[y_j,z_j]^2$ and
$\Frob_{\mathfrak t_k}^4$ map to $1$, since $\mathrm{Dih}_4$ has class two
and exponent four and its commutators lie in $\langle z\rangle$. So $\pi_M$
factors through $G_B$. Since $D_3(\mathrm{Dih}_4)=1$ (its Frattini subgroup
is $\langle z\rangle$, which is central of order two), $\pi_M$ kills
$D_3G_B$, and $M$ lies in the fixed field of $D_3G_B$, which is contained in
every $K\in\mathcal K$. On $D_2G_B$, $\pi_M$ is a homomorphism to
$\langle z\rangle$ that kills $D_3G_B$, hence factors through
$\operatorname{gr}_2G_B$; it agrees with $g\mapsto z^{\psi([g]_2)}$ on
squares and commutators, by the formulas above, and these classes span
$\operatorname{gr}_2G_B$.

(c) The prime $\mathfrak p$ is unramified in $M$. If it splits completely in
$F_0$, then the Frobenius of a prime $\mathfrak P_M$ of $M$ above a prime
$\mathfrak P$ of $F_0$ above $\mathfrak p$ lies in $\Gal(M/F_0)=\langle
z\rangle$, and it is $z$ exactly when $\beta$, a unit at $\mathfrak P$, is not
a square modulo $\mathfrak P$. In that case $f_M(\mathfrak p)=2$, and $2$
divides $f_K(\mathfrak p)$ because $M\subseteq K$
(Lemma~\ref{an:local-contribution}). The primes of $F_0$ above $\mathfrak p$
are conjugate under $\Gal(F_0/B)$, and $\sigma(\beta)/\beta$ is a square, so
the condition does not depend on $\mathfrak P$.

(d) If the Frobenius vector of $\mathfrak p$ is $0$, a Frobenius element
$g\in G_B$ of a prime of $B_G$ above $\mathfrak p$ lies in $D_2G_B$, and the
primes above $\mathfrak p$ give conjugate elements, whose classes in
$\operatorname{gr}_2G_B$ agree because $[G_B,D_2G_B]\subseteq D_3G_B$. By (b),
$\pi_M(g)=z^{\psi(\overline{\Frob}^{(2)}_{\mathfrak p})}$, and $\pi_M(g)$ is
the Frobenius of a prime of $M$ above $\mathfrak p$, which is $z$ exactly when
the hypothesis of (c) holds.
\end{proof}

\begin{lemma}\label{dh:beta-exists}
For each of the $27$ forms of Lemma~\ref{dh:27}, with the data
$(r_1,r_2,\mathbf r)$ of \texttt{d4all27.json}, there is an element $\beta$
satisfying the hypotheses of Lemma~\ref{dh:detection}. For the fifteen
census forms, the forms numbered $0,5,9,17,22,14,3,7,4,15,6,20,12,24,13$
(Definition~\fref{dh:census-basis}), and the seven forms of
Table~\ref{dh:W-table}, the elements used in this paper are given by their
coordinates in \texttt{d4fields15.json} and in \texttt{vd4\_data27.gp}; the
forms $7$, $17$, $20$ and $24$ belong to both lists, so these are $18$
distinct elements, checked $22$ times. For each of them, the ideal
factorization involves only primes above $2$, $3$ and $5$, and for every
$\sigma\ne1$ in $\Gal(F_0/B)$ the quotient $\sigma(\beta)/\beta$ is a square
and the sign $t_\sigma$ is the one required by (D2).
\end{lemma}

\begin{proof}[Proof (computer-assisted)]
The supplementary program \texttt{d4beta.gp} finds the elements: it computes
the $S$-units of $F_0$ modulo squares, the elements whose conjugates are
squares up to the element itself, and the linear map
$\beta\mapsto(t_\sigma)_\sigma$, and solves for the pattern of $\mathbf r$.
The programs \texttt{vd4.gp} (for the census forms) and \texttt{vplane4.gp}
(for the seven forms of Table~\ref{dh:W-table}) verify each of the $22$
(form, element) pairs in an independently built model of $F_0$: that $\beta$ is a unit
outside the primes above $S$ (by PARI's \texttt{idealfactor}), that the
quotients $\sigma(\beta)/\beta$ are squares with the stated signs
$t_\sigma$, and that the relative discriminant of $F_0(\sqrt\beta)/F_0$ has
norm divisible only by $2$, $3$ and $5$.
\end{proof}

\subsection{\texorpdfstring{The space $W''$}{The space W''}}\label{dh:W-section}
This subsection defines the space $W''$ of linear forms on
$\operatorname{gr}_2G_B$ spanned by four forms of dihedral type, its
subspaces $W\subset W'$ used for the intermediate results of
Section~\ref{lim:section}, and the field $\EW$ (Definition~\ref{dh:W}), and
describes the fifteen nonzero elements of $W''$
(Lemma~\ref{dh:W-structure}). The census (Subsection~\ref{dh:census-section})
decides which primes lie in $\ker W''$, and
Subsection~\ref{dh:EW-section} studies $\EW$.

\begin{table}[htbp]
\centering\small
\caption{The seven forms $\psi_o$ of dihedral type in $W''$, numbered as in
\texttt{d4all27.json}, with data $(r_{1,o},r_{2,o},\mathbf r_o)$, the vectors
written as strings of their coordinates as in Section~\ref{tw:tower}, and the
element $c_o\in B$ with $B_{c_o}=B(\sqrt{c_o})$ of Lemma~\ref{dh:Lfunctions}.
The forms satisfy $\psi_{24}=\psi_{10}+\psi_{23}$,
$\psi_{20}=\psi_{10}+\psi_{19}$ and $\psi_7=\psi_{17}+\psi_{19}$.}
\label{dh:W-table}
\begin{tabular}{rllll}
\toprule
$o$ & $r_{1,o}$ & $r_{2,o}$ & $\mathbf r_o$ & $c_o$\\
\midrule
$10$ & $00010011$ & $01001100$ & $(1,1)$ & $213033204-13722675\sqrt{241}$\\
$23$ & $01011111$ & $10001011$ & $(0,1)$ & $155-10\sqrt{241}$\\
$24$ & $01011111$ & $10011000$ & $(0,1)$ & $(-378557+24385\sqrt{241})/2$\\
$19$ & $01001100$ & $10000010$ & $(1,1)$ & $-326+21\sqrt{241}$\\
$20$ & $01001100$ & $10010001$ & $(1,1)$ & $(47-3\sqrt{241})/2$\\
$17$ & $01000110$ & $10001000$ & $(1,1)$ & $-31+2\sqrt{241}$\\
$7$ & $00001010$ & $11000100$ & $(1,1)$ & $-3433342+221161\sqrt{241}$\\
\bottomrule
\end{tabular}
\end{table}

\begin{definition}\label{dh:W}
Let $\psi_o$, $o\in\{10,23,24,19,20,17,7\}$, be the forms of
Table~\ref{dh:W-table}, with data $(r_{1,o},r_{2,o},\mathbf r_o)$ and
associated maps $\phi_o(v)=(r_{1,o}\cdot v,r_{2,o}\cdot v)$. Let
$F_{0,o}=B(\sqrt{\alpha_{r_{1,o}}},\sqrt{\alpha_{r_{2,o}}})$, let
$\beta_o\in F_{0,o}$ be the element of Lemma~\ref{dh:beta-exists}, and let
$M_o=F_{0,o}(\sqrt{\beta_o})$ be the dihedral field of
Lemma~\ref{dh:detection}. Put
\[
 W=\langle\psi_{10},\psi_{23}\rangle\subset
 W'=\langle\psi_{10},\psi_{23},\psi_{19}\rangle\subset
 W''=\langle\psi_{10},\psi_{23},\psi_{19},\psi_{17}\rangle.
\]
Let $W''^\perp\subset\operatorname{gr}_2G_B$ be the common kernel of the forms
in $W''$, let $N_{W''}\subseteq D_2G_B$ be its preimage, and let $\EW$ be the
fixed field of $N_{W''}$ in $B_G$. A prime $\mathfrak p\notin S$ of $B$ with
Frobenius vector $0$ \emph{lies in} $\ker W''$ if
$\overline{\Frob}^{(2)}_{\mathfrak p}\in W''^\perp$
(Lemma~\ref{dh:detection}(d)), that is, if
$\omega(\overline{\Frob}^{(2)}_{\mathfrak p})=0$ for every $\omega\in W''$;
likewise for $\ker W$ and $\ker W'$.
\end{definition}

\begin{lemma}\label{dh:W-structure}
The space $W''$ has dimension $4$. Its $15$ nonzero elements are the seven
forms of Table~\ref{dh:W-table}, which are of dihedral type, and the eight
forms $\psi_a+\psi_b$ for the pairs
\[
 (a,b)=(19,23),\ (19,24),\ (7,10),\ (17,10),\ (17,24),\ (7,24),\ (7,23),\
 (17,23),
\]
whose polar forms have rank $4$; for each of these pairs the four vectors
$r_{1,a},r_{2,a},r_{1,b},r_{2,b}$ span a space of dimension $4$, and
$c_a/c_b$ is not a square in $B$. Every subspace of the dual of
$\operatorname{gr}_2G_B$ that properly contains $W''$ contains a form whose
polar form has rank at least $6$.
\end{lemma}

\begin{proof}[Proof (computer-assisted)]
The supplementary program \texttt{plane4.py} computes the $15$
elements and their polar ranks from \texttt{d4all27.json}, and checks the last
statement by running over all $2^{15}$ linear forms; \texttt{vplane4.gp}
verifies the seven dihedral fields (Lemma~\ref{dh:beta-exists}). The rank of
the four vectors and the fields $B(\sqrt{c_a},\sqrt{c_b})$, of degree $8$
over $\Q$, are checked by \path{deg8/export8.gp}.
\end{proof}

\subsection{The census}\label{dh:census-section}
This subsection determines, for the primes of $B$ of degree one with
Frobenius vector $0$ and norm at most $4\cdot10^{13}$, whether they lie in
$\ker W''$ and whether fifteen fixed dihedral fields detect them
(Proposition~\ref{dh:census}). The detected primes in $\ker W''$ receive the
floor $2$ in Definition~\ref{dh:floors}, and the census determines the
residue degrees in $\EW$ of the primes of small norm
(Lemma~\ref{dh:floors-valid}).

\begin{definition}\label{dh:census-basis}
The \emph{census forms} are $\psi^{\rm c}_i=\psi_{o_i}$, $1\le i\le15$, where
\[
 (o_1,\dots,o_{15})=(0,5,9,17,22,14,3,7,4,15,6,20,12,24,13)
\]
is the order of the file \texttt{d4fields15.json}, which records their data
and elements $\beta^{\rm c}_1,\dots,\beta^{\rm c}_{15}$
(Lemma~\ref{dh:beta-exists}); let $F^{\rm c}_{0,i}$ be the field $F_0$ of
Lemma~\ref{dh:detection} for $\psi^{\rm c}_i$. (The superscript distinguishes
these from $\beta_o$ and $F_{0,o}$ of Definition~\ref{dh:W}, indexed by the form
number $o$.) For a prime $\mathfrak p$ of $B$ of degree one with
Frobenius vector $0$, its \emph{pattern} is
\[
 \operatorname{pat}(\mathfrak p)
 =\bigl(\psi^{\rm c}_i(\overline{\Frob}^{(2)}_{\mathfrak p})\bigr)_{i\le15}
 \in\Ftwo^{15},
\]
and $\mathfrak p$ is \emph{detected} if $\operatorname{pat}(\mathfrak p)\ne0$.
\end{definition}

\begin{lemma}\label{dh:vector-zero}
Let $p\ge7$ be a prime with $(241/p)=1$, let $r^2\equiv241\pmod p$, and let
$\mathfrak p=(p,\sqrt{241}-r)$. Then $\mathfrak p$ has Frobenius vector $0$
exactly when $p\equiv1$ or $49\pmod{120}$ and the images of
$\alpha_1,\alpha_2,\alpha_4,\alpha_6$ under $\sqrt{241}\mapsto r$ are squares
modulo $p$. In that case the conjugate prime $(p,\sqrt{241}+r)$ also has
Frobenius vector $0$.
\end{lemma}

\begin{proof}
The residue field of $\mathfrak p$ is $\mathbb F_p$, and by
Definition~\ref{gf:frobenius-vector} and Euler's criterion $\mathfrak p$ has
Frobenius vector $0$ exactly when $\alpha_0,\dots,\alpha_7$ are squares
modulo $\mathfrak p$; they are units at $\mathfrak p$ since $p\ge7$ and
$p\ne241$. By \eqref{tw:alpha}, $\alpha_0=-1$, $\alpha_2\alpha_3=2$,
$\alpha_4\alpha_5=-3$, $\alpha_6\alpha_7=-5$ and
$\mathrm N_{B/\Q}(\alpha_1)=-1$. If $\mathfrak p$ has vector $0$, then $-1$, $2$,
$-3$ and $-5$ are squares modulo $p$, so $p\equiv1\pmod8$, $p\equiv1\pmod3$
and $p\equiv\pm1\pmod5$, that is, $p\equiv1$ or $49\pmod{120}$. Conversely,
for such $p$ the numbers $-1,2,3,5$ are squares modulo $p$, and the
relations show that $\alpha_0,\alpha_3,\alpha_5,\alpha_7$ are squares modulo
$\mathfrak p$ when $\alpha_2,\alpha_4,\alpha_6$ are. By \eqref{tw:alpha},
the nontrivial automorphism of $B$ maps $\alpha_1,\alpha_2,\alpha_4,\alpha_6$ to
$-\alpha_1^{-1}$, $-\alpha_3=-2\alpha_2^{-1}$, $\alpha_5=-3\alpha_4^{-1}$ and
$\alpha_7=-5\alpha_6^{-1}$, which are squares modulo $\mathfrak p$ exactly
when $\alpha_1,\alpha_2,\alpha_4,\alpha_6$ are; so the conjugate prime also
has vector $0$.
\end{proof}

\begin{lemma}\label{dh:pattern-lemma}
Let $\mathfrak p$ be a prime of $B$ of degree one with Frobenius vector $0$,
and $p=\mathrm N\mathfrak p$, with $\operatorname{pat}(\mathfrak p)$ as in
Definition~\ref{dh:census-basis}.
\begin{enumerate}
\item[(a)] $\operatorname{pat}(\mathfrak p)_i=1$ exactly when $\beta^{\rm c}_i$ is a
 nonsquare modulo a prime of $F^{\rm c}_{0,i}$ above $\mathfrak p$; these primes have
 residue field $\mathbb F_p$.
\item[(b)] $\operatorname{pat}(\mathfrak p)=0$ exactly when
 $\overline{\Frob}^{(2)}_{\mathfrak p}=0$, that is, when the Frobenius
 elements of the primes of $B_G$ above $\mathfrak p$ lie in $D_3G_B$.
\item[(c)] If $\psi=\sum_{i\in I}\psi^{\rm c}_i$ for a set $I$ of indices,
 then $\psi(\overline{\Frob}^{(2)}_{\mathfrak p})=\sum_{i\in I}
 \operatorname{pat}(\mathfrak p)_i$.
\item[(d)] If $\mathfrak p$ is detected, then $f_K(\mathfrak p)$ is even for
 every $K\in\mathcal K$.
\item[(e)] $\psi_{17}=\psi^{\rm c}_4$,
 $\psi_{19}=\psi^{\rm c}_4+\psi^{\rm c}_8$,
 $\psi_{10}=\psi^{\rm c}_4+\psi^{\rm c}_8+\psi^{\rm c}_{12}$ and
 $\psi_{23}=\psi^{\rm c}_4+\psi^{\rm c}_8+\psi^{\rm c}_{12}+\psi^{\rm c}_{14}$;
 the index sets are encoded, with bit $i-1$ for $\psi^{\rm c}_i$, by the
 masks $0x8$, $0x88$, $0x888$ and $0x2888$ of the file
 \texttt{wmasks4.json}. Hence $\mathfrak p$ lies in $\ker W''$ exactly when
 $\operatorname{pat}(\mathfrak p)$ has even weight on each of these four
 index sets.
\end{enumerate}
\end{lemma}

\begin{proof}
(a) The prime $\mathfrak p$ splits completely in $E\supseteq F^{\rm c}_{0,i}$, so the
primes of $F^{\rm c}_{0,i}$ above it have residue field $\mathbb F_p$, and the
statement is Lemma~\ref{dh:detection}(d) for $\psi^{\rm c}_i$ and $\beta^{\rm c}_i$.
(b) The supplementary program \texttt{d4fields.py} checks that the fifteen
census forms have rank $15$, so they form a basis of the dual of
$\operatorname{gr}_2G_B$, which has dimension $15$; hence
$\operatorname{pat}(\mathfrak p)=0$ exactly when
$\overline{\Frob}^{(2)}_{\mathfrak p}=0$. (c) is linearity. (d) Some
$\operatorname{pat}(\mathfrak p)_i$ is $1$, and Lemma~\ref{dh:detection}(c)
applies to $\psi^{\rm c}_i$ by (a). (e) We have $o_4=17$, $o_8=7$, $o_{12}=20$
and $o_{14}=24$, and the relations of Table~\ref{dh:W-table} give
$\psi_{19}=\psi_{17}+\psi_7$, $\psi_{10}=\psi_{19}+\psi_{20}$ and
$\psi_{23}=\psi_{10}+\psi_{24}$. The last statement follows from (c), since
$\psi_{10},\psi_{23},\psi_{19},\psi_{17}$ span $W''$.
\end{proof}

\begin{proposition}[The census]\label{dh:census}
The primes of $B$ of degree one with Frobenius vector $0$ and norm at most
$4\cdot10^{13}$, sorted by whether they lie in $\ker W''$
(Definition~\ref{dh:W}) and are detected (Definition~\ref{dh:census-basis}),
are as follows:
\begin{center}\small
\begin{tabular}{lrrr}
\toprule
range of norms & vector $0$ & in $\ker W''$ & in $\ker W''$, detected\\
\midrule
$65809\le\mathrm N\mathfrak p\le1.2\cdot10^7$ & $2924$ & $136$ & $136$\\
$(1.2\cdot10^7,10^{13}]$ & $1351669408$ & $84431399$ & $84395189$\\
$(10^{13},4\cdot10^{13}]$ & $3807510358$ & $237942878$ & $237834284$\\
\bottomrule
\end{tabular}
\end{center}
The smallest norm of a prime of degree one with vector $0$ is $65809$. The
undetected primes are those in $\ker W''$ that are not detected; all
other primes of the table are detected or lie outside $\ker W''$. The norms of
the detected primes in $\ker W''$ are recorded exactly for norms at most
$1.2\cdot10^7$, and above that in the geometric bins
$[X(1+2^{-16})^j,X(1+2^{-16})^{j+1})$, $X=1.2\cdot10^7$ or $10^{13}$, as counts
per bin (files \texttt{census/w4\_1e13.binsBW4} and
\texttt{w4\_4e13.binsBW4}).
\end{proposition}

\begin{proof}[Proof (computer-assisted)]
The census runs over the primes $\mathfrak p=(p,\sqrt{241}-r)$ of degree
one, $r^2\equiv241\pmod p$, and tests vector $0$ by
Lemma~\ref{dh:vector-zero}. Square roots modulo $p$ of the images of the
Kummer generators $\alpha_0,\dots,\alpha_7$ under $\sqrt{241}\mapsto r$ give
an embedding of each $F^{\rm c}_{0,i}$ into $\mathbb F_p$; each square root is
checked by squaring, and $\beta^{\rm c}_i$ is evaluated from its coordinates, so that
$\operatorname{pat}(\mathfrak p)$ is computed with Legendre symbols modulo
$p$ (Lemma~\ref{dh:pattern-lemma}(a)). By Lemma~\ref{dh:beta-exists},
$\beta^{\rm c}_i$ is a unit at every prime above $p\ge7$, so no Legendre symbol is
$0$; every run reports this.
Membership in $\ker W''$ is decided by Lemma~\ref{dh:pattern-lemma}(e). For
norms at most $1.2\cdot10^7$, the supplementary program
\path{dihedral_ceiling3.py} classifies the primes in Python with
the functions of \texttt{rootcheck.py}. Above $1.2\cdot10^7$, the C program
\texttt{census\_kv4.c} produced the bins; it computes, eight primes
at a time in Montgomery arithmetic with the AVX-512 IFMA instructions, the
vector-$0$ test and the Legendre symbols of the $\beta^{\rm c}_i$: first for the four
census fields whose bits determine $\psi_{10}$, $\psi_{23}$, $\psi_{19}$ and
$\psi_{17}$, and then, for the primes on which $\psi_{10}$ and $\psi_{23}$
vanish, for the other census fields until one detects. In the bins a norm
may be assigned to a neighboring bin by floating-point rounding; the bound of
Lemma~\fref{ce:values} allows one bin of slack on each side. The program and
its data were checked as follows; the earlier programs named here are those
of the intermediate certificates.
\begin{enumerate}
\item[(1)] On the vector-$0$ primes of norm $6\cdot10^4$ to $1.2\cdot10^7$, all
 C programs reproduce the Python classification prime by prime and field by
 field, and \texttt{census\_kv4.c}'s bins for $\ker W$, $\ker W'$ and
 $\ker W''$ (Definition~\ref{dh:W}) equal an independent classification by
 \texttt{rootcheck.py}.
\item[(2)] Its bins for $\ker W$ on $(1.2\cdot10^7,10^{13}]$ and
 $(10^{13},4\cdot10^{13}]$ are byte-identical to those of the earlier programs
 \texttt{census\_d3k.c} and \texttt{census\_kw.c}, which use scalar
 arithmetic and a different structure. On sample ranges near $10^{13}$ and
 $5\cdot10^{13}$, \texttt{census\_kv4.c} and \texttt{census\_kv.c} agree in
 both their checking and their fast modes. In the fast mode about one
 vector-$0$ prime in sixteen, and in the checking mode every one, is also
 recomputed by the scalar code, and the outcomes must agree; all sixteen
 square roots of the Kummer elements are checked by squaring, which
 re-proves vector $0$.
\item[(3)] Independent reviews recounted bins chosen at random with their own
 programs and found exact agreement, and a second implementation in PARI/GP
 matched patterns, bins and undetected primes on eight ranges up to
 $10^{12}$.
\end{enumerate}
The replay program reruns (1) and the comparisons of (2) on sample ranges,
and checks the stored bins by their SHA-256 hashes.
\end{proof}

\subsection{\texorpdfstring{The field $E_{W''}$ and its $L$-functions}{The field E\_W'' and its L-functions}}
\label{dh:EW-section}
This subsection shows that the field $\EW$ of Definition~\ref{dh:W} is
Galois over $B$, of degree $8192$ over $\Q$, and contained in every
$K\in\mathcal K$; it factors $\zeta_{\EW}$ into $\zeta_E$ and $576$
$L$-functions (Lemma~\ref{dh:EW}), identifies these as Hecke $L$-functions
of quadratic characters with known conductors and gamma factors
(Lemmas~\ref{dh:Lfunctions} and~\ref{dh:gamma}), whose values
Section~\ref{lv:section} encloses, and computes residue degrees in $\EW/B$
(Lemma~\ref{dh:residue}).

\begin{lemma}\label{dh:EW}
\begin{enumerate}
\item[(a)] $\EW$ is Galois over $B$ and contained in every $K\in\mathcal K$.
 It contains $E$ and the seven dihedral fields $M_o$ of
 Definition~\ref{dh:W}. Its group $G_{W''}=\Gal(\EW/B)$ is a central
 extension of $\Gal(E/B)=\Ftwo^8$ by $Z_{W''}=D_2G_B/N_{W''}\cong\Ftwo^4$, and
 $[\EW:\Q]=8192$.
\item[(b)] For $\omega\in W''\setminus\{0\}$, regarded as a character of
 $Z_{W''}=\operatorname{gr}_2G_B/W''^\perp$, let $r(\omega)$ be the rank of
 its polar form $\operatorname{pol}_\omega$. The irreducible complex
 representations of $G_{W''}$ with central character $\omega$ have dimension
 $2^{r(\omega)/2}$, and there are $256/2^{r(\omega)}$ of them. Those with
 trivial central character are the $256$ characters $\chi_w$.
\item[(c)] For $\omega=\psi_o$ of dihedral type, they are the $64$
 representations $\rho_o\otimes\chi_w$, where $\rho_o$ is the
 two-dimensional irreducible representation of
 $\Gal(M_o/B)\cong\mathrm{Dih}_4$, inflated to $G_{W''}$, and $w$ runs over a
 complement of the span of $r_{1,o},r_{2,o}$.
\item[(d)] For $\omega=\psi_a+\psi_b$ as in Lemma~\ref{dh:W-structure}, they
 are the $16$ representations $\rho_a\otimes\rho_b\otimes\chi_w$, with $w$
 over a complement of the span of $r_{1,a},r_{2,a},r_{1,b},r_{2,b}$.
\item[(e)] With $L(s,\pi)$ the Artin $L$-function over $B$,
 \[
  \zeta_{\EW}(s)=\zeta_E(s)\prod_{o}\prod_{w}L(s,\rho_o\otimes\chi_w)^2
  \prod_{(a,b)}\prod_{w}L(s,\rho_a\otimes\rho_b\otimes\chi_w)^4,
 \]
 the products over the seven forms and the eight pairs, and over the
 complements of (c) and (d). Hence, for real $s>1$,
 \begin{equation}\label{dh:YW}
  Y_{W''}(s):=\frac{\log\zeta_{\EW}(s)}{8192}
  =\frac{Y_E(s)}{16}+\frac1{4096}\sum_{o,w}\log L(s,\rho_o\otimes\chi_w)
  +\frac1{2048}\sum_{(a,b),w}\log L(s,\rho_a\otimes\rho_b\otimes\chi_w),
 \end{equation}
 where $Y_E(s)=\log\zeta_E(s)/512$ as in Proposition~\ref{an:genus-value}.
\end{enumerate}
\end{lemma}

\begin{proof}
(a) Since $[G_B,D_2G_B]\subseteq D_3G_B\subseteq N_{W''}$, every subgroup
between $D_3G_B$ and $D_2G_B$ is normal in $G_B$, and $D_2G_B/N_{W''}$ is
central in $G_B/N_{W''}$; so $\EW$ is Galois over $B$ and $G_{W''}$ is a central
extension of $G_B/D_2G_B=\Gal(E/B)$ by $Z_{W''}$. The latter is
$\operatorname{gr}_2G_B/W''^\perp$, which is dual to $W''$, of dimension
four. Hence $[\EW:B]=2^{12}$ and $[\EW:\Q]=8192$. Since $N_{W''}\supseteq
D_3G_B\supseteq D_4G_B$, $\EW$ lies in every $K\in\mathcal K$. By
Lemma~\ref{dh:detection}(b), the kernel of $\phi_{M_o}$ contains the elements
of $D_2G_B$ whose class is killed by $\psi_o$, hence contains $N_{W''}$, so
$M_o\subseteq\EW$.

(b) The commutator of elements of $G_{W''}$ with elementary images $v,w$ is
the class of $[v,w]$ in $Z_{W''}$, and the square of an element with image $v$
is the class of $v^{[2]}$. Let $\omega\ne0$, and let $H=G_{W''}/\ker\omega$.
The representations with central character $\omega$ factor through $H$, a
central extension of $\Ftwo^8$ by $\langle z\rangle\cong C_2$ in which
elements with images $v,w$ have commutator $z^{\operatorname{pol}_\omega(v,w)}$.
The center $C_H$ of $H$ is the preimage of the radical of $\operatorname{pol}_\omega$,
of index $2^{r(\omega)}$ in $H$. Let $\pi$ be irreducible with $\pi(z)=-1$.
If $g\notin C_H$, there is $h$ with $h^{-1}gh=gz$, so
$\tr\pi(g)=-\tr\pi(g)=0$; on $C_H$, $\pi$ acts by scalars, so
$|\tr\pi|=\dim\pi$ there. The orthogonality relation
$\sum_{g\in H}|\tr\pi(g)|^2=|H|$ gives $(\dim\pi)^2=[H:C_H]=2^{r(\omega)}$.
The irreducible representations of $H$ with $z\mapsto-1$ have
$\sum(\dim\pi)^2=|H|-|H/\langle z\rangle|=2^8$, so there are
$2^8/2^{r(\omega)}$ of them. Those with trivial central character are the
representations of $\Gal(E/B)=\Ftwo^8$.

(c) The representation $\rho_o\otimes\chi_w$ has central character
$\psi_o$, by Lemma~\ref{dh:detection}(b), and dimension $2=2^{r(\psi_o)/2}$,
so it is irreducible. Its character vanishes outside the preimage of the
radical $\ker\phi_o$ of $\operatorname{pol}_{\psi_o}$
(Lemma~\ref{dh:dihedral-polar}), and is $\pm2\chi_w$ on it. Hence
$\rho_o\otimes\chi_w\cong\rho_o\otimes\chi_{w'}$ exactly when
$\chi_{w-w'}$ is trivial on $\ker\phi_o$, that is, when $w-w'$ lies in the
span of $r_{1,o},r_{2,o}$. This gives $256/4=64$ classes, all of them by (b).

(d) Likewise $\rho_a\otimes\rho_b\otimes\chi_w$ has central character
$\psi_a+\psi_b$ and dimension $4=2^{r(\psi_a+\psi_b)/2}$, so it is irreducible. The radical of
the polar form, of codimension four, contains $\ker\phi_a\cap\ker\phi_b$, which
has codimension four by Lemma~\ref{dh:W-structure}, so they are equal, and
the twists are distinct modulo the span of the four vectors.

(e) The Artin formalism gives $\zeta_{\EW}=\prod_\pi L(s,\pi)^{\dim\pi}$ over
the irreducible representations of $G_{W''}$, and $\prod_wL(s,\chi_w)=
\zeta_E(s)$ (Proposition~\ref{gf:factorization}). By Lemma~\ref{dh:W-structure} every nonzero $\omega\in W''$ is
of one of the two kinds in (c) and (d), which gives the product; the count
$256+7\cdot64\cdot4+8\cdot16\cdot16=4096=|G_{W''}|$ is a check. Taking
logarithms and dividing by $8192=16\cdot512$ gives \eqref{dh:YW}.
\end{proof}

\begin{remark}\label{dh:EW-sub}
The proof of Lemma~\ref{dh:EW} uses only that every nonzero element of
$W''$ is a form of dihedral type or a sum $\psi_a+\psi_b$ as in
Lemma~\ref{dh:W-structure}. It therefore applies to every subspace
$U\subseteq W''$, for instance to $W$ and $W'$, with $E_U$ the fixed field of
the preimage in $D_2G_B$ of the common kernel of the forms in $U$: the field
$E_U$ is Galois over $B$, contained in every $K\in\mathcal K$, of degree
$512\cdot2^{\dim U}$ over $\Q$, and $\zeta_{E_U}$ is $\zeta_E$ times the
squares of the $64$ functions $L(s,\rho_o\otimes\chi_w)$ of each form
$\psi_o\in U$ and the fourth powers of the $16$ functions
$L(s,\rho_a\otimes\rho_b\otimes\chi_w)$ of each sum $\psi_a+\psi_b\in U$.
\end{remark}

\begin{lemma}\label{dh:Lfunctions}
Let $\rho_o$ and $\chi_w$ be as in Lemma~\ref{dh:EW}.
\begin{enumerate}
\item[(a)] Let $\psi_o$ be a form of Table~\ref{dh:W-table}. Let
 $\tau_o\in\Gal(F_{0,o}/B)$ be one of the two nonzero elements other than
 $\mathbf r_o$, and let $B_{c_o}=B(\sqrt{c_o})$ be its fixed field in
 $F_{0,o}$; Table~\ref{dh:W-table} records the element $c_o$ used. Let
 $u_o\in F_{0,o}^\times$ satisfy $\tau_o(\beta_o)/\beta_o=u_o^2$ and
 $u_o\tau_o(u_o)=1$ (Lemma~\ref{dh:detection}, (D1) and (D2)), with
 $u_o\ne-1$ (replace $u_o$ by $-u_o$ if necessary). Then
 $\gamma_o=(1+u_o)^2\beta_o$ lies in $B_{c_o}$, and for every $w$,
 \[
  L(s,\rho_o\otimes\chi_w)=\frac{\zeta_{M_{o,w}}(s)}{\zeta_{B_{c_o}}(s)},\qquad
  M_{o,w}=B_{c_o}\bigl(\sqrt{\gamma_o\alpha_w}\bigr).
 \]
\item[(b)] For a pair $(a,b)$ of Lemma~\ref{dh:W-structure}, put
 $B_{ab}=B(\sqrt{c_a},\sqrt{c_b})$. Then for every $w$
 \[
  L(s,\rho_a\otimes\rho_b\otimes\chi_w)
  =\frac{\zeta_{B_{ab,w}}(s)}{\zeta_{B_{ab}}(s)},
  \qquad B_{ab,w}=B_{ab}\bigl(\sqrt{\gamma_a\gamma_b\alpha_w}\bigr).
 \]
\item[(c)] Each of these quotients is the Hecke $L$-function $L(s)$ of a
 nontrivial quadratic character of $A=B_{c_o}$, respectively $A=B_{ab}$, and
 it is entire. Let $A'$ be the quadratic extension $M_{o,w}$, respectively
 $B_{ab,w}$, of $A$, and $Q=|\Delta_{A'}|/|\Delta_A|$. Then the function
 $\Lambda(s)=Q^{s/2}\Gamma_\R(s)^{r_1(A')-r_1(A)}\Gamma_\C(s)^{r_2(A')-r_2(A)}L(s)$
 satisfies $\Lambda(1-s)=\Lambda(s)$, where $(r_1(\cdot),r_2(\cdot))$ is the
 signature.
\end{enumerate}
\end{lemma}

\begin{proof}
(a) Fix $o$ and write $F_0$, $\beta$, $\tau$, $u$, $\gamma$ and $c$ for
$F_{0,o}$, $\beta_o$, $\tau_o$, $u_o$, $\gamma_o$ and $c_o$. Let
$\tilde\tau\in\Gal(M_o/B)$ lift $\tau$ with $\tilde\tau(\sqrt\beta)=u\sqrt\beta$.
Then $\tilde\tau^2(\sqrt\beta)=\tau(u)u\sqrt\beta=\sqrt\beta$, so $\tilde\tau$
has order two, and
$\tilde\tau((1+u)\sqrt\beta)=(1+\tau(u))u\sqrt\beta=(1+u)\sqrt\beta$. So the
nonzero element $(1+u)\sqrt\beta$ is fixed by $\tilde\tau$; its square
$\gamma$ lies in $F_0$ and is fixed by $\tau$, so it lies in $B_c$; and
$\sqrt\gamma\notin B_c$, since $\sqrt\beta\notin F_0$. The fixed field of
$\langle\tilde\tau\rangle$ in $M_o$ is therefore $B_c(\sqrt\gamma)$. Let
$H_o\subset G_{W''}$ be the fixed group of $B_c$, of index two, and let
$\xi_o$ be the quadratic character of $H_o$ that cuts out $B_c(\sqrt\gamma)$.
The image of $H_o$ in $\Gal(M_o/B)\cong\mathrm{Dih}_4$ is
$\langle\tilde\tau,z\rangle$, and $\xi_o$ factors through it, trivial on
$\tilde\tau$ and nontrivial on $z$. Hence $\Ind_{H_o}^{G_{W''}}\xi_o$ is
inflated from the representation of $\mathrm{Dih}_4$ induced from this
character of $\langle\tilde\tau,z\rangle$, which has dimension two and is
nontrivial on the center, so it is the two-dimensional irreducible
representation: $\Ind_{H_o}\xi_o=\rho_o$. Twisting,
$\rho_o\otimes\chi_w=\Ind_{H_o}(\xi_o\cdot\chi_w|_{H_o})$, and
$\xi_o\cdot\chi_w|_{H_o}$ is the quadratic character of $H_o$ that cuts out
$M_{o,w}$. It is nontrivial, since $\rho_o\otimes\chi_w$ is irreducible. By
the invariance of Artin $L$-functions under induction,
$L(s,\rho_o\otimes\chi_w)=L(s,\xi_o\chi_w|_{H_o})$, which is
$\zeta_{M_{o,w}}(s)/\zeta_{B_c}(s)$.

(b) Write $\rho_a=\Ind_{H_a}\xi_a$ and $\rho_b=\Ind_{H_b}\xi_b$ as in (a). Since
$c_a/c_b$ is not a square, $H_a\ne H_b$, so $H_aH_b=G_{W''}$, and Mackey's
formula \cite[Section~7.3]{Serre1977} gives
$\operatorname{Res}_{H_a}\Ind_{H_b}\xi_b=\Ind_{H_a\cap H_b}^{H_a}\xi_b$. Hence
\[
 \rho_a\otimes\rho_b\otimes\chi_w
 =\Ind_{H_a}\bigl(\xi_a\otimes\operatorname{Res}_{H_a}(\rho_b\otimes\chi_w)\bigr)
 =\Ind_{H_a\cap H_b}\bigl(\xi_a\xi_b\chi_w\bigr),
\]
and $H_a\cap H_b$ is the fixed group of $B_{ab}$; the character
$\xi_a\xi_b\chi_w$ cuts out $B_{ab,w}$, and it is nontrivial since the
representation is irreducible.

(c) A nontrivial Hecke character has an entire $L$-function
\cite{Hecke1920,Tate1967}. The completed Dedekind zeta functions satisfy
$\Lambda_A(1-s)=\Lambda_A(s)$ (Section~\ref{gf:section}), and the quotient of
those of $A'$ and $A$ is the stated $\Lambda$.
\end{proof}

\begin{lemma}\label{dh:gamma}
The conductors and gamma factors of the $576$ functions of
Lemma~\ref{dh:Lfunctions} are as follows. For $o=10,23,24,19,17,7$, all $64$
twists have $r_1(M_{o,w})=r_1(B_{c_o})$, so their gamma factor is
$\Gamma_\C(s)^2$; for $o=20$, $B_{c_o}$ is totally real and
$r_1(M_{o,w})=6$ for $32$ twists and $2$ for the other $32$, so the gamma
factor is $\Gamma_\R(s)^3\Gamma_\R(s+1)$ or $\Gamma_\R(s)\Gamma_\R(s+1)^3$. For
the eight pairs, $B_{ab}$ and all $B_{ab,w}$ are totally complex, so the
gamma factor is $\Gamma_\C(s)^4$. The conductors lie in the following
ranges:
\begin{center}\small
\begin{tabular}{lll}
\toprule
functions & degree over $\Q$ & conductors\\
\midrule
$o=10$ & $4$ & $4.18\cdot10^8$ to $2.68\cdot10^{10}$\\
$o=23,24$ & $4$ & $1.67\cdot10^9$ to $2.68\cdot10^{10}$\\
$o=19,7$ & $4$ & $1.67\cdot10^8$ to $2.68\cdot10^{11}$\\
$o=20$ & $4$ & $3.35\cdot10^8$ to $1.34\cdot10^{11}$\\
$o=17$ & $4$ & $6.69\cdot10^8$ to $2.68\cdot10^{11}$\\
pairs $(19,23)$, $(19,24)$ & $8$ & $6.995\cdot10^{17}$ and $1.791\cdot10^{20}$\\
the other six pairs & $8$ & $2.798\cdot10^{18}$ or $4.477\cdot10^{19}$, and $7.163\cdot10^{20}$\\
\bottomrule
\end{tabular}
\end{center}
\end{lemma}

\begin{proof}[Proof (computer-assisted)]
The twists of $\psi_{10},\psi_{23},\psi_{24}$ are treated by the
supplementary program \texttt{export.gp}, those of the other four forms by
the program \path{deg8/exportg.gp}, and the pairs by
\path{deg8/export8.gp}. These programs construct $B_{c_o}$, $\gamma_o$,
$M_{o,w}$, $B_{ab}$ and $B_{ab,w}$ in PARI/GP from the data of
Lemmas~\ref{dh:detection} and~\ref{dh:beta-exists}, and compute the
signatures, the discriminants and the Euler factors at small primes from the
maximal orders. For the twists of the seven forms the maximal orders come
from PARI's \texttt{nfinit}. For the pairs, whose fields have degree $16$,
the program computes orders maximal at $2,3,5,241$ and asserts that their
discriminants are, up to sign, products of powers of these primes, so that
the orders are maximal; it also asserts the signatures and that each local
quotient of Euler factors is a polynomial of degree at most $8$ in $p^{-s}$.
Wrong conductors or gamma factors would also be detected by the checks of
Subsection~\ref{lv:checks}.
\end{proof}

\begin{lemma}\label{dh:residue}
Let $\mathfrak p\notin S$ be a prime of $B$ with Frobenius vector $v$, and let
$f_{W''}(\mathfrak p)$ be its residue degree in the extension $\EW/B$ of
Lemma~\ref{dh:EW}, with the maps $\phi_o$ of Definition~\ref{dh:W}.
\begin{enumerate}
\item[(a)] If $v\ne0$, then $f_{W''}(\mathfrak p)=4$ if
 $\phi_o(v)=\mathbf r_o$ for one of the seven forms $\psi_o$ of
 Table~\ref{dh:W-table}, and $f_{W''}(\mathfrak p)=2$ otherwise.
\item[(b)] If $v=0$, then $f_{W''}(\mathfrak p)=2$ if $\mathfrak p$ does not lie
 in $\ker W''$, and $f_{W''}(\mathfrak p)=1$ if it does.
\end{enumerate}
\end{lemma}

\begin{proof}
Let $g\in G_{W''}$ be a Frobenius of a prime above $\mathfrak p$; its order is
$f_{W''}(\mathfrak p)$. (a) If $v\ne0$, then $g\notin Z_{W''}$ and $g^2\in
Z_{W''}$, and $\omega(g^2)=q_\omega(v)$ for $\omega\in W''$. So $g^2=1$
exactly when $q_\omega(v)=0$ for all $\omega\in W''$, that is, for the seven
dihedral forms, which span $W''$; and $q_{\psi_o}(v)=1$ exactly when
$\phi_o(v)=\mathbf r_o$. (b) If $v=0$, then $g\in Z_{W''}$, and $g\ne1$
exactly when $\omega(\overline{\Frob}^{(2)}_{\mathfrak p})\ne0$ for some
$\omega\in W''$.
\end{proof}

\subsection{Floors}\label{dh:floors-section}
This subsection defines the floors used in the proof of
Theorem~\ref{an:ceiling} (Definition~\ref{dh:floors}), shows that they form a
system of floors (Lemma~\ref{dh:floors-valid}), and expresses their floor
sum through $Y_{W''}$ and explicit corrections
(Lemmas~\ref{dh:floor-sum} and~\ref{dh:delta-sel}); Section~\ref{ce:section}
evaluates the result. For primes of degree one with vector $0$,
Lemma~\ref{dh:residue}(b) is read from the census
(Lemma~\ref{dh:pattern-lemma}(e)); the primes of degree two with vector $0$
and small norm are treated by the next lemma.

\begin{lemma}\label{dh:inert-v0}
There are exactly thirteen primes of $B$ of degree two, that is,
$p\mathcal O_B$ with $p$ inert in $B$, with Frobenius vector $0$ and norm
$p^2\le1.2\cdot10^7$, namely $p=409$, $769$, $1129$, $1249$, $1489$, $1609$,
$1801$, $2281$, $2521$, $2689$, $3001$, $3049$ and $3361$. None of them lies
in $\ker W''$, so $f_{W''}(p\mathcal O_B)=2$ for each of them.
\end{lemma}

\begin{proof}[Proof (computer-assisted)]
For such a prime $\mathfrak p=p\mathcal O_B$, the Frobenius in
$\Gal(M_o/B)$ is central by Lemma~\ref{dh:detection}(b), so it is $1$ or $z$,
and accordingly the Euler factor at $p$ of $L(s,\rho_o\otimes\chi_0)$ is
$(1-p^{-2s})^{-2}$ or $(1+p^{-2s})^{-2}$; by Lemma~\ref{dh:detection}(b) it
is $z$ exactly when $\psi_o(\overline{\Frob}^{(2)}_{\mathfrak p})=1$. The
supplementary program \texttt{inertv0\_4.py} lists the inert primes of
vector $0$ with $p^2\le1.2\cdot10^7$, computes these Euler factors with PARI
from the prime decompositions of $p$ in $M_{o,0}$ and $B_{c_o}$
(Lemma~\ref{dh:Lfunctions}(a)), and finds that none of the thirteen primes
lies in $\ker W''$. Lemma~\ref{dh:residue}(b) then gives
$f_{W''}(p\mathcal O_B)=2$.
\end{proof}

\begin{definition}\label{dh:floors}
Let $\mathcal P_2$ be the set of the primes of $B$ of degree one with vector $0$
and norm at most $4\cdot10^{13}$ that lie in $\ker W''$ and are detected; they
are counted in the last column of Proposition~\fref{dh:census}. For a free
prime $\mathfrak p$ of $B$ put $f^0_{\mathfrak p}=4$ if
$\mathfrak p\in\mathcal P_4$ (Definition~\ref{an:P4}),
$f^0_{\mathfrak p}=2$ if $\mathfrak p\in\mathcal P_2$, and
$f^0_{\mathfrak p}=f_{W''}(\mathfrak p)$ otherwise.
\end{definition}

\begin{lemma}\label{dh:floors-valid}
The numbers $f^0_{\mathfrak p}$ form a system of floors
(Definition~\ref{an:floor-def}), and $f^0_{\mathfrak p}\ge f_{W''}(\mathfrak
p)$ for every free $\mathfrak p$.
\end{lemma}

\begin{proof}
(Fl1) The residue degree $f_{W''}(\mathfrak p)$ divides $f_K(\mathfrak p)$,
since $\EW\subseteq K$ (Lemmas~\ref{dh:EW}(a) and~\ref{an:local-contribution}).
For $\mathfrak p\in\mathcal P_4$, $4$ divides $f_K(\mathfrak p)$ by
Lemma~\ref{an:census-degree}; for $\mathfrak p\in\mathcal P_2$, $2$ divides
$f_K(\mathfrak p)$ by Lemma~\ref{dh:pattern-lemma}(d). Since
$f_{W''}(\mathfrak p)\le4$ always and $f_{W''}(\mathfrak p)=1$ on
$\mathcal P_2$, the floors are at least $f_{W''}(\mathfrak p)$. (Fl2) For
$\mathrm N\mathfrak p>1.2\cdot10^7>41^4$ it is clear. For
$\mathrm N\mathfrak p\le1.2\cdot10^7$ the supplementary program
\path{dihedral_ceiling3.py} computes $f^0_{\mathfrak p}$ by
Lemma~\ref{dh:residue}, using Proposition~\ref{dh:census} for the primes of
degree one with vector $0$ and Lemma~\ref{dh:inert-v0} for those of degree
two, and checks $\mathrm N\mathfrak p^{f^0_{\mathfrak p}}\ge41^4$; in
particular every free prime of norm less than $1681$ has floor $4$.
\end{proof}

\begin{lemma}\label{dh:floor-sum}
For real $\sigma>1$, the floor sum of the floors of
Definition~\ref{dh:floors} is
\[
 Y_{\rm fl}(\sigma)=Y_{W''}(\sigma)-\Delta_{\rm sel}(\sigma)-Z_{\rm sel}(\sigma)
 -\Delta_4(\sigma)-\sum_{\mathfrak p\in\mathcal P_2}T_\sigma(\mathrm N\mathfrak p),
\]
where $T_\sigma(N)=g_\sigma(N)/2-g_\sigma(N^2)/4$,
\[
 \Delta_4(\sigma)=\sum\Bigl(\frac{g_\sigma(\mathrm N\mathfrak p^2)}4
 -\frac{g_\sigma(\mathrm N\mathfrak p^4)}8\Bigr),
\]
the sum over the $\mathfrak p\in\mathcal P_4$ with $f_{W''}(\mathfrak p)=2$,
and $\Delta_{\rm sel}(\sigma)$ is the sum over the selected primes
$\mathfrak p$ of the difference between the term of $\mathfrak p$ in
\eqref{an:local-sum} for $M=\EW$ and for $M=K$, with the types of
Theorem~\ref{tw:field-family}.
\end{lemma}

\begin{proof}
By \eqref{an:local-sum} for $M=\EW$, $Y_{W''}(\sigma)$ is the sum over all
primes of $B$ of the terms $g_\sigma(\mathrm N\mathfrak p^f)/(2ef)$ with the
type $(e,f)$ of $\mathfrak p$ in $\EW/B$. For the selected primes, the term
for $\EW$ is the term for $K$ plus a summand of $\Delta_{\rm sel}(\sigma)$, and
the terms for $K$ add up to $Z_{\rm sel}(\sigma)$ by
Lemma~\ref{an:selected-lemma}. A free prime has $e=1$ in $\EW$, and its
floor differs from $f_{W''}(\mathfrak p)$ only on $\mathcal P_4$, from $2$ to
$4$, and on $\mathcal P_2$, from $1$ to $2$; these changes lower its term by the
summands of $\Delta_4(\sigma)$ and by $T_\sigma(\mathrm N\mathfrak p)$. The
sets $\mathcal P_4$ and $\mathcal P_2$ are disjoint, since the primes of
$\mathcal P_4$ have nonzero vectors.
\end{proof}

\begin{lemma}\label{dh:delta-sel}
For real $\sigma>1$,
\[
 \Delta_{\rm sel}(\sigma)=Y^{(2)}_{W''}(\sigma)-\frac{g_\sigma(16)}{32}
 +\sum_{k=1}^3\Bigl(\frac{g_\sigma(\mathrm N\mathfrak t_k^{f_k})}{2f_k}
 -\frac{g_\sigma(\mathrm N\mathfrak t_k^4)}8\Bigr),
\]
where $f_k=f_{W''}(\mathfrak t_k)\in\{2,4\}$ is given by
Lemma~\ref{dh:residue}(a), and $Y^{(2)}_{W''}(\sigma)$ is the right side of
\eqref{dh:YW} at $s=\sigma$ with $\zeta_E$ and each $L$-function replaced by
the product of its Euler factors at the primes of $B$ above $2$. The
selected primes above $3$ and $5$ contribute $0$ to $\Delta_{\rm sel}$.
\end{lemma}

\begin{proof}
By \eqref{an:local-sum} for $M=\EW$, restricted to the primes above $2$, the
terms of $\mathfrak p_1,\mathfrak p_2$ for $M=\EW$ add up to the logarithm
of the product of the Euler factors of $\zeta_{\EW}$ at the primes above $2$,
divided by $8192$; by the factorization of Lemma~\ref{dh:EW}(e), which holds
Euler factor by Euler factor, this is $Y^{(2)}_{W''}(\sigma)$. Their terms for
$K$, of type $(8,4)$ and norm $2$, are $g_\sigma(16)/64$ each. Above $3$ and
$5$ the types in $E$ (Corollary~\ref{gf:types}) and in $K$ are both $(2,2)$,
so the type in $\EW$ is $(2,2)$ too, since $E\subseteq\EW\subseteq K$ and
ramification indices and residue degrees divide in towers; these primes
contribute nothing. The capped primes have nonzero Frobenius vectors
(Table~\ref{tw:vector-table}), so their residue degree $f_k$ in $\EW$ is $2$
or $4$ by Lemma~\ref{dh:residue}(a), with $e=1$; their type in $K$ is
$(1,4)$.
\end{proof}

The supplementary program \path{dihedral_ceiling3.py} evaluates
$\Delta_{\rm sel}$ by Lemma~\ref{dh:delta-sel}, reading the Euler factors at
$2$ of all the $L$-functions in \eqref{dh:YW}; it also checks that the Euler
factors at $3$ and $5$ give the types $(2,2)$. For orientation, at
$\sigma=1001/1000$ it gives $Y_E=0.0853580483$, $Y_{W''}=0.050492599$ (radius
$3.4\cdot10^{-10}$), $\Delta_{\rm sel}=0.0079701070$ and
$\Delta_4=0.0005116777$.
